% AUTORIDAD DE FUENTE: Exposé IX §9.4, líneas francesas 2435--2576;
% índices de escaneo 438--443, folios 427--432. Las correcciones D/C
% de Hom, Ext y la referencia final se conservan mediante \sisi.

%%%%%%%%%%  scan idx 438  |  folio 427  %%%%%%%%%%

\medskip
\noindent
9.4. Sea $S$ un trait henseliano de punto genérico $\eta$, sea $n$ un
entero $>0$, pongamos $\Lambda=\mathbb Z/n\mathbb Z$, y sean
$\underline C$ la categoría de $\Lambda$-módulos sobre $S$ (para el
sitio fppf de $S$; véase VIII 0.2) y $\underline C'$ la categoría de
$\Lambda$-módulos sobre $\eta$. Sean $P,R,Q$ objetos de $\underline C$,
y sean $E,F$ extensiones de $Q$ por $R$ y de $R$ por $P$,
respectivamente (9.3.2). Supongamos que $Q$ es un $\Lambda_S$-módulo
localmente libre de tipo finito y que $P$ está representado por un grupo
finito de tipo multiplicativo sobre $S$ (SGA 3 IX), de modo que es de la
forma
\begin{equation*}
 P=D(Q^*),
\end{equation*}
donde $D=\underline{\Hom}(-,\mathbb G_{mS})$ denota la dualidad de
Cartier. Denotemos por $P_\eta,R_\eta,Q_\eta,E_\eta,F_\eta$ las
restricciones de $P,R,Q,E,F$ a $\eta$, y supongamos dada una extensión
mixta $X_\eta$ de $E_\eta$ por $F_\eta$. En estas condiciones definiremos
canónicamente un emparejamiento
\begin{equation}\tag{9.4.1}
 c(X_\eta):Q\otimes_\Lambda Q^*\longrightarrow\Lambda_S
 \qquad(\Lambda=\mathbb Z/n\mathbb Z).
\end{equation}

Supongamos primero que $S$ es \underline{estrictamente local}. Entonces
$Q$ y $Q^*$ son constantes, digamos $Q=Q_{0S}$ y $Q^*=Q^*_{0S}$, y dar
(9.4.1) equivale a dar un elemento
\begin{equation}\tag{9.4.2}
 c(X_\eta)\in\check Q_0\otimes\check Q_0^*,
\end{equation}
donde el signo $\check{\ }$ denota el dual ordinario
$\Hom(-,\Lambda)$. Aplicamos las consideraciones de 9.3 al funtor de
restricción $r:\underline C\longrightarrow\underline C'$. Observemos
primero que existe una extensión mixta de $E$ por $F$. En efecto, la
obstrucción a su existencia pertenece a $\Ext^2(Q,P)$ (9.3.8(c)), y
basta demostrar
\begin{equation*}
 \Ext^2(Q,P)=0.
\end{equation*}

Esto es un caso particular del lema siguiente.

%%%%%%%%%%  scan idx 439  |  folio 428  %%%%%%%%%%

\medskip
\noindent
\underline{Lema 9.4.3. Bajo las hipótesis precedentes, con $S$
estrictamente local, se tienen:}

\smallskip
\noindent
a) \underline{isomorfismos canónicos}
\begin{equation*}\tag{a)}
 \Ext^i_\Lambda(Q,P)\simeq H^i(S,\check Q\otimes P),
 \qquad
 \Ext^i_\Lambda(Q_\eta,P_\eta)
 \simeq H^i(\eta,\check Q\otimes P);
\end{equation*}

\noindent
b) \underline{un isomorfismo canónico}
\begin{equation}\tag{b)}
 \check Q\otimes P\simeq D(Q\otimes Q^*);
\end{equation}

\noindent
c) \underline{para todo grupo finito constante $H=H_{0S}$ sobre $S$,}
\begin{equation}\tag{c)}
 H^i(S,D(H))=H^i(\eta,D(H)_\eta)=0
 \qquad\underline{\text{para }}i\geq2;
\end{equation}

\noindent
d) \underline{isomorfismos canónicos}
\begin{equation}\tag{d)}
\begin{aligned}
 H^1(S,D(H))&\simeq\mathbb G_m(S)\otimes\check H_0
               =V^*\otimes\check H_0,\\
 H^1(\eta,D(H)_\eta)&\simeq\mathbb G_m(\eta)\otimes\check H_0
               =K^*\otimes\check H_0,
\end{aligned}
\end{equation}
\underline{donde en (d) el signo $\check{\ }$ denota
$\Hom(-,\mathbb Q/\mathbb Z)$.}

Las fórmulas de (a) y (b) se deducen inmediatamente de que $Q$ es un
$\Lambda_S$-módulo localmente libre de tipo finito. Para demostrar (c),
nos reducimos al caso $H_0=\mathbb Z/m\mathbb Z$, es decir,
$D(H)=\mu_m$. La sucesión de cohomología asociada a la sucesión exacta
de Kummer
\begin{equation*}
 0\longrightarrow\mu_m\longrightarrow\mathbb G_m
 \xrightarrow{\ m\cdot\id\ }\mathbb G_m
 \longrightarrow0
\end{equation*}
nos reduce a demostrar
\begin{equation*}
 H^i(S,\mathbb G_m)=H^i(\eta,\mathbb G_m)=0
 \qquad\text{para }i\geq1.
\end{equation*}
Estos grupos son los $H^i_{\mathrm{\acute et}}$ calculados en el sitio
étale [6, Teorema 11.7], y por tanto se anulan para $S$, pues $S$ es
estrictamente local. Su anulación para $\eta$ también es bien conocida
[6, Teorema 1.1].

%%%%%%%%%%  scan idx 440  |  folio 429  %%%%%%%%%%

Queda por describir los isomorfismos de (d). Sobre una base arbitraria
$U$ existe un homomorfismo canónico
\begin{equation}\tag{9.4.3.1}
 \mathbb G_m(U)\otimes\check H_0
 \longrightarrow H^1(U,D(H_{0U})).
\end{equation}
Si $H_0$ es anulado por un entero $m>0$, este homomorfismo se describe
mediante el morfismo de borde
$\mathbb G_m(U)\xrightarrow{\partial}H^1(U,\mu_m)$ asociado a la
sucesión de Kummer: es el compuesto
\begin{equation*}
\begin{aligned}
 \mathbb G_m(U)\otimes\check H_0
 &\longrightarrow H^1(U,\mu_m)\otimes\check H_0\\
 &\longrightarrow H^1(U,\mu_m\otimes\check H_0)
  \longrightarrow H^1(U,D(H_{0U})).
\end{aligned}
\end{equation*}
La segunda flecha se deduce del homomorfismo evidente
$\check H_0\longrightarrow
\Hom(\mu_m,\mu_m\otimes\check H_0)$, y la tercera es un caso particular
de (b), con $m$ en lugar de $n$. Dejamos al lector comprobar que
(9.4.3.1) es independiente de la elección del entero $m$ tal que
$mH_0=0$. Afirmamos también que (9.4.3.1) es siempre inyectivo y que su
cociente se anula si $\Pic(U)=0$. De hecho, dicho cociente es siempre
isomorfo a $\Hom(H_0,\Pic(U))$, pero esto no importa aquí. En efecto, nos
reducimos al caso $H_0=\mathbb Z/n\mathbb Z$, y la afirmación se sigue
inmediatamente de la sucesión de cohomología de Kummer. Esto demuestra
(d) y da inmediatamente:

\medskip
\noindent
\underline{Corolario 9.4.4. Con la notación de 9.4.3, existe un
isomorfismo canónico}

\noindent
e)
\begin{equation*}\tag{e)}
 H^1(\eta,D(H)_\eta)/\operatorname{Im}H^1(S,D(H))
 \simeq\check H_0,
\end{equation*}
\underline{y, combinando} (a) \underline{y} (b), \underline{un isomorfismo canónico}
\begin{equation}\tag{9.4.5}
 \Ext^1(Q_\eta,P_\eta)/
 \operatorname{Im}\Ext^1(Q,P)
 \simeq\check Q_0\otimes\check Q_0^*.
\end{equation}

%%%%%%%%%%  scan idx 441  |  folio 430  %%%%%%%%%%

Volvamos a la definición de (9.4.2) cuando $S$ es estrictamente local.
La relación $\Ext^2(Q,P)=0$ permite definir el elemento canónico
(9.3.13), tomando como $Y'$ la extensión mixta dada $X_\eta$ de
$E_\eta$ por $F_\eta$. La fórmula (9.4.5) identifica este elemento
$c(X_\eta)$ con un elemento de $\check Q_0\otimes\check Q_0^*$, lo que
completa la definición de (9.4.2). Las consideraciones generales de 9.3
muestran también que este elemento es la obstrucción a prolongar
$X_\eta$ a una extensión mixta $X$ de $E$ por $F$.

\medskip
\noindent
9.4.6. La definición de (9.4.1) cuando $S$ es henseliano se reduce
inmediatamente al caso estrictamente local por descenso de Galois desde
el henselizado estricto de $S$.

\medskip
\noindent
\underline{Proposición 9.4.7. Sean $P,Q,R,E,F$ como antes, salvo que el
trait $S$ no tiene por qué ser henseliano. Consideremos el funtor}
\begin{equation}\tag{9.4.7.1}
 X\longmapsto X_\eta
\end{equation}
\underline{de la categoría de extensiones mixtas de $E$ por $F$ a la
categoría de extensiones mixtas de $E_\eta$ por $F_\eta$. Para toda
extensión mixta $X_\eta$ de $E_\eta$ por $F_\eta$, consideremos el
emparejamiento}
\begin{equation}\tag{9.4.7.2}
 c_0(X_\eta):Q_0\otimes Q_0^*\longrightarrow\Lambda_s
\end{equation}
\underline{inducido por el emparejamiento de tipo (9.4.1) definido
sobre el henselizado $S^h$ de} $S$. \underline{Entonces el funtor (9.4.7.1) es
plenamente fiel, y una extensión mixta $X_\eta$ de $E_\eta$ por
$F_\eta$ pertenece a su imagen esencial si y solo si el emparejamiento
(9.4.7.2) es nulo.}

%%%%%%%%%%  scan idx 442  |  folio 431  %%%%%%%%%%

9.4.7.3. Este enunciado sigue teniendo sentido y siendo válido si se
supone solamente que $S$ es un esquema conexo, regular, noetheriano y de
dimensión $1$. En vez del único emparejamiento (9.4.7.2), hay que
considerar un emparejamiento
$Q_s\otimes Q_s^*\longrightarrow\Lambda_s$ para cada punto cerrado $s$
de $S$. Será más cómodo demostrar 9.4.7 en esta forma más general.

\smallskip
\noindent
a) Para toda extensión mixta $X$, la aplicación
\begin{equation*}
 \End(X)\longrightarrow\End(X_\eta)
\end{equation*}
es biyectiva. Sus dos miembros son, respectivamente, las secciones sobre
$S$ y $\eta$ de
$\underline{\Hom}(\sisi{P,Q}{Q,P})=D(Q\otimes Q^*)$; por tanto, la
afirmación se sigue de que $S$ es normal y $D(Q\otimes Q^*)$ es finito
sobre $S$. De aquí se deduce inmediatamente que, si $X$ e $Y$ son dos
extensiones mixtas que ya se sabe isomorfas, entonces
\begin{equation*}
 \Hom(X,Y)\longrightarrow\Hom(X_\eta,Y_\eta)
\end{equation*}
es biyectiva.

\smallskip
\noindent
b) Si $X$ e $Y$ son dos extensiones mixtas, todo isomorfismo
$X_\eta\xrightarrow{\sim}Y_\eta$ procede de un isomorfismo
$X\longrightarrow Y$. Por (a), la cuestión es local para la topología
étale, por lo que podemos reducirnos al caso en que $S$ es estrictamente
local. En ese caso, 9.4.3 muestra que la aplicación
\begin{equation*}
 \Ext^1_\Lambda(\sisi{P,Q}{Q,P})
 \longrightarrow
 \Ext^1_\Lambda(\sisi{P_\eta,Q_\eta}{Q_\eta,P_\eta})
\end{equation*}
se identifica con
\begin{equation*}
 V^*\otimes\check H_0\longrightarrow K^*\otimes\check H_0,
 \qquad H_0=Q_0\otimes Q_0^*.
\end{equation*}
Esta última es inyectiva, pues la sucesión exacta
$0\longrightarrow V^*\longrightarrow K^*\longrightarrow\mathbb Z
\longrightarrow0$ es escindida. Por la discusión de 9.3, esto implica
que, si $X_\eta$ e $Y_\eta$ son isomorfas, también lo son $X$ e $Y$; la
conclusión se sigue entonces de (a).

Así, (a) y (b) muestran que $X\longmapsto X_\eta$ es
\underline{plenamente} fiel.

%%%%%%%%%%  scan idx 443  |  folio 432  %%%%%%%%%%

\smallskip
\noindent
c) Sea $X_\eta$ una extensión mixta sobre $\eta$. Demostremos que se
prolonga a $S$ si y solo si los emparejamientos
$(\sisi{9.4.6.2}{9.4.7.2})$ son nulos. También aquí, la unicidad ya
demostrada muestra que la cuestión es local para la topología étale, lo
que nos reduce de nuevo al caso estrictamente local. La conclusión vale
entonces por construcción, como se observó inmediatamente antes de
9.4.6.

Esto concluye la demostración de 9.4.7.
